\section{What systems of linear equations describe}

A system of linear equations is not only a collection of equations to solve. It is a collection of conditions that must be satisfied at the same time. Each equation gives one condition. A solution is an assignment of values to the unknowns that satisfies all conditions simultaneously.

This point is important. When we solve a system, we are not solving the equations one after another as separate problems. We are looking for values that make the whole system true.

\begin{examplebox}
Consider the system
\[
\begin{cases}
x+y=5,\\
2x-y=1.
\end{cases}
\]
The pair \((2,3)\) satisfies the first equation because
\[
2+3=5.
\]
But it does not satisfy the second equation, because
\[
2\cdot2-3=1.
\]
Actually this equality is true, so \((2,3)\) satisfies both equations. Therefore \((2,3)\) is a solution of the system.

The pair \((1,4)\) satisfies the first equation, since \(1+4=5\), but it does not satisfy the second equation, since \(2\cdot1-4=-2\). Therefore \((1,4)\) is not a solution of the system.
\end{examplebox}

The example should be read slowly. A solution of a system must pass every test in the system. Passing one equation is not enough.

\subsection*{Unknowns, coefficients, and constants}

A linear equation in the unknowns \(x_1,x_2,\ldots,x_n\) has the form
\[
a_1x_1+a_2x_2+\cdots+a_nx_n=b.
\]
The numbers \(a_1,a_2,\ldots,a_n\) are coefficients. The number \(b\) is the constant term or right-hand side.

A system of linear equations is made of several such equations:
\[
\begin{cases}
a_{11}x_1+a_{12}x_2+\cdots+a_{1n}x_n=b_1,\\
a_{21}x_1+a_{22}x_2+\cdots+a_{2n}x_n=b_2,\\
\hspace{2cm}\vdots\\
a_{m1}x_1+a_{m2}x_2+\cdots+a_{mn}x_n=b_m.
\end{cases}
\]
Here \(m\) is the number of equations and \(n\) is the number of unknowns.

\begin{remarkbox}
The coefficients describe how the unknowns enter the equations. The right-hand side describes what values the linear combinations must equal. A system compares these two pieces of information.
\end{remarkbox}

\subsection*{Solution sets}

The solution set of a system is the set of all solutions. Sometimes this set contains one element. Sometimes it is empty. Sometimes it contains infinitely many elements.

These are the three basic possibilities for a linear system:
\begin{itemize}
\item exactly one solution,
\item no solution,
\item infinitely many solutions.
\end{itemize}

The goal of solving a system is not only to find numbers. The goal is to determine which of these cases occurs and to describe the solution set accurately.

\subsection*{Matrix form}

Systems of linear equations are naturally connected with matrices. The coefficients form a matrix, the unknowns form a column vector, and the right-hand side forms another column vector.

For example, the system
\[
\begin{cases}
2x+y=5,\\
x-y=1
\end{cases}
\]
can be written as
\[
\begin{pmatrix}
2&1\\
1&-1
\end{pmatrix}
\begin{pmatrix}
x\\y
\end{pmatrix}
=
\begin{pmatrix}
5\\1
\end{pmatrix}.
\]
Using the notation
\[
A=\begin{pmatrix}2&1\\1&-1\end{pmatrix},
\qquad
x=\begin{pmatrix}x\\y\end{pmatrix},
\qquad
b=\begin{pmatrix}5\\1\end{pmatrix},
\]
we write the system compactly as
\[
Ax=b.
\]

This notation is not merely shorter. It reveals the structure of the problem. The matrix \(A\) transforms the unknown vector \(x\) into the right-hand side vector \(b\). Solving the system means finding which input vector \(x\), if any, is sent to \(b\).

\subsection*{The augmented matrix}

When solving a system, it is often convenient to write the coefficients and the right-hand side together in an augmented matrix:
\[
\left[
\begin{array}{cc|c}
2&1&5\\
1&-1&1
\end{array}
\right].
\]
The vertical line reminds us that the last column is not another column of coefficients. It is the right-hand side.

The augmented matrix is a compact record of the system. Row operations on this matrix correspond to controlled transformations of the equations.

\begin{summarybox}
When reading a system of linear equations, ask:
\begin{itemize}
\item What are the unknowns?
\item What are the equations, and what condition does each one impose?
\item What are the coefficients and right-hand sides?
\item What would it mean for one assignment to satisfy all equations at once?
\item Can the system be written as \(Ax=b\)?
\item What does the augmented matrix record?
\end{itemize}
\end{summarybox}

A system of equations should first be read as a structured mathematical object. Only then should we choose a method for solving it.